\begin{figure}[tb]
\centering
\begin{tikzpicture}
\begin{axis}[axis main,width=.78\linewidth,height=6.3cm,axis equal image,
 xlabel={$\eta$},ylabel={$\chi$},xmin=0,xmax=5.1,ymin=-2.1,ymax=2.1]
 \addplot[loss curve,domain=0:360,samples=121]
 ({1+0*cos(x)},{0+0*sin(x)});
 \addplot[loss curve,MidnightBlue!65,domain=0:360,samples=121]
 ({1.25+.75*cos(x)},{.75*sin(x)});
 \addplot[loss curve,MidnightBlue!40,domain=0:360,samples=121]
 ({1.75+1.436*cos(x)},{1.436*sin(x)});
 \addplot[voltage curve,domain=0:360,samples=121]
 ({3.15+1.05*cos(x)*cos(18)-.62*sin(x)*sin(18)},
  {.15+1.05*cos(x)*sin(18)+.62*sin(x)*cos(18)});
 \addplot[only marks,mark=*,ForestGreen] coordinates {(1,0)};
 \node[annotation,anchor=west] at (axis cs:1.05,.05){$J=2$};
 \node[annotation] at (axis cs:3.55,-1.25){voltage conic};
\end{axis}
\end{tikzpicture}
\caption{Projective loss levels are circles with moving centers. The
unconstrained optimum is the degenerate point $(1,0)$; voltage forces the first
contact with its transformed conic.}
\label{fig:projective}
\end{figure}
